\section{Budgets and shapes}\label{sec:structure}

\subsection{Attention against transport}
Adaptive transport networks also strengthen links that are used \citep{tero2010rules}. Under a steady load they form
trees; fluctuating loads keep loops \citep{katifori2010damage,corson2010fluctuations}. E6 writes both as one update,
$W\leftarrow(1-\eta)W+\eta\,\mathcal N[U^p]$, and swaps two ingredients (\cref{fig:e6flow}): the budget $\mathcal N$ (per
agent, as attention, or one shared pool, as material) and the usage $U$ (the info score, or the squared Kirchhoff current
when the source broadcasts). The per-agent/closeness cell is the attention model with tempered selection ($T=1/p$); the
shared/flow cell has Corson's optimum as its steady-load fixed point ($\gamma=1/p-1$). Runs use a $7\times7$ lattice,
$p\in[0.4,20]$, $\eta=0.05$, 3000 steps, three repeats, with constant, Markov and random sources.

\figw{e6-flow}{\linewidth}{The shared update of E6; the four cells are the four choices of usage (box 2) and budget
(box 4).}{e6flow}

\figpair{e6-loops-constant}{constant source}{e6-loops-random}{random source}{Independent loops of the backbone against
concentration $p$ (maximum 36). Shared-budget points that collapse onto part of the graph are omitted.}{e6loops}

\begin{itemize}
  \item \textbf{Trees need self-reinforcement.} The flow reward reinforces itself; closeness does so only through
    $\kappa$. At $\kappa=0$ equally close parents tie forever and every loop survives; trees appear past $p\approx1/\kappa$,
    where the symmetric state loses stability ($T=\kappa$).
  \item \textbf{A fluctuating source keeps loops only with the flow reward.} With closeness, a random source merely delays
    the transition; past it the network locks into some tree whatever the source does.
  \item \textbf{The budget decides robustness and cost.} A per-agent budget forces every agent to keep a link: strong
    concentration gives trees, never a collapse, with dissipation at most about 4.5 times the uniform lattice. A shared
    budget puts the tree transition exactly at Corson's $\gamma=1$ but collapses at strong concentration, with dissipation
    growing by orders of magnitude.
\end{itemize}

\subsection{Paths in the hierarchy morphospace}
Following one experimental variable at a time through $(T_h,F,O)$ shows where each experiment moves the shape
(\cref{fig:paths}). Commitment is a one-dimensional descent from the hierarchy corner to about $(0.25,0.25,0.25)$ as $T$
passes $\Thalf$; noise lowers feedforwardness first; and a hard selection cut climbs toward the same region under a random
source as under a constant one, so that structure comes from selection alone.

\figw{paths}{\linewidth}{Paths in the hierarchy morphospace: E8 over commitment at 80 agents, one line per noise level
(left), and E5 over commitment with a source moving every 100 rounds, one line per $\kappa$ (right). Hollow markers start
each path.}{paths}

\subsection{Shape is not function}
Under a moving source, runs below the lock-in edge sit at the same coordinates as true hierarchies, and cost the same: a
hub on a stale agent is as tree-like as a hub on the source. Position in the morphospace must be read together with
tracking.
